% ============================================================
% pretty-charts 字体层：衬线搭配（宋体 + Times New Roman）
% 用法：\input{.../references/style/tikz/fonts-serif.tex}
%
% 为什么单独成文件：示例此前各自写 \setmainfont{Times New Roman}
% + \setCJKmainfont{SimSun}，这两个字体只存在于 Windows，同一份示例在
% Linux/macOS 上无法编译（CI 也就无从编译）。这里按"字体是否存在"逐级回退：
%   Windows  → Times New Roman + SimSun（与原行为完全一致）
%   Linux    → TeX Gyre Termes + Noto Serif CJK SC（apt install fonts-noto-cjk-extra）
%   兜底     → TeX Gyre Termes + FandolSong（TeX Live 自带）
%   再兜底   → 只发警告、不设字体，保证编译不因缺字体而失败（中文可能显示异常）
% \IfFontExistsTF 由 fontspec 提供（preamble.tex 已加载 fontspec）。
% ============================================================
\IfFontExistsTF{Times New Roman}
  {\setmainfont{Times New Roman}}
  {\setmainfont{TeX Gyre Termes}}          % TeX Live 自带，Times 的开源等价物
\IfFontExistsTF{SimSun}
  {\setCJKmainfont{SimSun}}                % 中文宋体（Windows）
  {\IfFontExistsTF{Noto Serif CJK SC}
     {\setCJKmainfont{Noto Serif CJK SC}}   % Linux（fonts-noto-cjk-extra）
     {\IfFontExistsTF{FandolSong}
        {\setCJKmainfont{FandolSong}}        % 兜底：TeX Live 自带，无需系统字体包
        {\PackageWarning{pretty-charts}{未找到中文衬线字体（已尝试 SimSun / %
          Noto Serif CJK SC / FandolSong）；中文可能显示异常。%
          请安装 fonts-noto-cjk-extra，或改用 fonts-sans.tex}}}}
